\section{SoC Integration and Memory Map Concept}

\subsection{Why a CPU core needs a SoC wrapper}

A CPU alone only executes instructions. A usable system needs memory and peripherals. The SoC wrapper connects the CPU to:

\begin{itemize}
\item instruction memory;
\item data RAM;
\item UART;
\item timer/CLINT;
\item syscon/control registers;
\item debug blocks;
\item FPGA top-level pins.
\end{itemize}
\subsection{Memory-mapped peripherals}

The CPU uses normal load/store instructions. The SoC address decoder decides what hardware block responds.

\begin{lstlisting}[language={}]
CPU address -> address decoder -> RAM or UART or timer or syscon
\end{lstlisting}

This means software can write a UART transmit register with a normal store instruction.

\subsection{Important wrappers}

\begin{longtable}{|>{\RaggedRight\arraybackslash}p{0.38\textwidth}|>{\RaggedRight\arraybackslash}p{0.52\textwidth}|}
\hline
\textbf{File} & \textbf{Role} \\
\hline
\endfirsthead
\hline
\textbf{File} & \textbf{Role} \\
\hline
\endhead
\texttt{rtl/soc.v} & Basic simulation SoC around \texttt{cpu\_core.v}. \\
\hline
\texttt{rtl/soc\_pipe.v} & SoC wrapper for pipelined CPU experiments. \\
\hline
\texttt{rtl/soc\_mmu.v} & SoC wrapper for MMU experiments. \\
\hline
\texttt{rtl/soc\_fpga.v}, \texttt{rtl/soc\_uart\_fpga.v}, \texttt{rtl/soc\_fpga\_mmu.v} & FPGA-oriented system wrappers. \\
\hline
\texttt{rtl/soc\_rtos.v}, \texttt{rtl/soc\_rtos\_fpga.v} & RTOS-oriented systems with timer/UART integration. \\
\hline
\texttt{rtl/fpga\_top*.v} & Board-level top modules. \\
\hline
\end{longtable}

\bigskip
